\documentclass[11pt,a4paper]{article}
\usepackage[margin=2.5cm]{geometry}
\usepackage{amsmath}
\usepackage{booktabs}
\usepackage{listings}
\usepackage{xcolor}
\usepackage{hyperref}
\usepackage{graphicx}
\usepackage{parskip}

\hypersetup{colorlinks=true, linkcolor=blue, urlcolor=blue}

\lstset{
  basicstyle=\ttfamily\small,
  keywordstyle=\color{blue},
  commentstyle=\color{gray},
  stringstyle=\color{red},
  frame=single,
  numbers=left,
  numberstyle=\tiny,
  breaklines=true,
  language=C
}

\title{\textbf{Progress Report 2}\\
\large GPULAB: High Performance Computing with Graphic Cards\\
Project Topic 1-2: CT Volume Reconstruction}

\author{Sanan Akther, Divyansh Sharma, Joseph Maher Labib Estawro}
\date{July 15, 2026}

\begin{document}
\maketitle
\tableofcontents
\newpage

% ─────────────────────────────────────────────────────────────────────────────
\section{Progress Summary}

Since Progress Report 1, all four reconstruction modes have been fully implemented,
debugged, and validated. The CPU C port, GPU buffer mode, GPU image-object mode, and
GPU optimized mode all produce numerically correct output with Mean Squared Error
(MSE) below $2 \times 10^{-8}$ relative to the CPU reference on the $256^3$ dataset.

This report covers the C CPU implementation (Sections~\ref{sec:cpu}--\ref{sec:challenges}),
the three GPU OpenCL modes (Section~\ref{sec:gpu}), preliminary performance results
(Section~\ref{sec:perf}), and a correctness validation (Section~\ref{sec:validation}).
Remaining work (512$^3$ dataset, final optimization analysis, convergence plots) is
deferred to the final report.

% ─────────────────────────────────────────────────────────────────────────────
\section{CPU C Implementation}
\label{sec:cpu}

\subsection{Forward Projection (\texttt{fp\_cpu})}

The forward projection ray-marches through the 3D volume for each detector pixel
and angle. The outer loop is parallelized over $N_\theta \times W$ work items using
OpenMP \texttt{collapse(2)}, giving $75 \times 512 = 38{,}400$ independent tasks:

\begin{lstlisting}
#pragma omp parallel for collapse(2) schedule(static)
for (int ip = 0; ip < np; ip++) {
    for (int iu = 0; iu < W; iu++) {
        /* inner loop: ray march through Ny=512 samples */
        for (int iv = 0; iv < H; iv++) { ... }
    }
}
\end{lstlisting}

Per-sample trilinear interpolation uses a single unsigned bounds check (avoids
branching) and incremental ray stepping (avoids one multiply per sample):
\begin{lstlisting}
if ((unsigned)x0 < (unsigned)(Nxz-1) &&
    (unsigned)y0 < (unsigned)(Ny-1)  &&
    (unsigned)z0 < (unsigned)(Nxz-1)) {
    /* 8-tap trilinear */
}
wx += dox; wy += doy; wz += doz;  /* incremental step */
\end{lstlisting}

Result: $\approx 5.8$\,s/epoch (was $\approx 28$\,s in the naive implementation).

\subsection{Backprojection (\texttt{bp\_cpu})}

Backprojection is parallelized over $(ix, iy)$ voxel pairs with
\texttt{collapse(2)}. Each pair owns a unique z-strip of the output volume so no
atomics or reduction buffers are needed. A precomputed cos/sin LUT eliminates
per-voxel transcendental calls:

\begin{lstlisting}
#pragma omp parallel for collapse(2) schedule(static)
for (int ix = 0; ix < Nxz; ix++) {
    for (int iy = 0; iy < Nxz; iy++) {
        float *strip = volume + VOL_IDX(Nxz-1-ix, Nxz-1-iy, 0, Nxz, Ny);
        for (int ip = 0; ip < np; ip++) {
            float ca = ca_lut[ip], sa = sa_lut[ip];
            /* bilinear lookup into projection slice */
        }
        /* scale strip in-place */
    }
}
\end{lstlisting}

Result: $\approx 1.3$\,s/epoch. Total CPU: $\approx 7.3$\,s/epoch.

\subsection{MLEM Reconstruction Loop}

The iterative update follows the MLEM rule from the Python reference:
\[
v^{(k+1)} = v^{(k)} \cdot \frac{\mathcal{B}\!\left(\dfrac{p_0}{\mathcal{F}(v^{(k)})}\right)}{\mathcal{B}(\mathbf{1})}
\]

A critical subtlety: the Python \texttt{bp\_func} applies the cone-weight
correction \emph{inside} each backprojection call. Consequently, the cone weight
must be applied to the ratio $p_0 / \mathcal{F}(v)$ \emph{before} backprojection
each epoch, and $\mathcal{B}(\mathbf{1})$ must be computed from cone-weighted ones.
Applying cone weight once to the measured data $p_0$ (as a preprocessing step) is
incorrect and produces divergence.

% ─────────────────────────────────────────────────────────────────────────────
\section{Implementation Challenges}
\label{sec:challenges}

\subsection{Projection Data Layout}

The Python reference transposes and flips the projection data inside \texttt{bp\_func}:
\begin{verbatim}
proj[:, ::-1, :].transpose(0,2,1) / voxelSize
\end{verbatim}
This converts the HDF5 layout $[N_\theta][H][W]$ to $[N_\theta][W][H]$ with a
flipped H-axis, then scales by $1/\delta_v$. This preprocessing must be applied to
both the ratio and the all-ones buffer before every backprojection call. A dedicated
\texttt{preprocess\_proj} OpenCL kernel handles this for GPU modes.

\subsection{Forward Projection Boundary Condition}

The OpenCL image3D sampler with \texttt{CLK\_ADDRESS\_CLAMP\_TO\_EDGE} extends
edge voxels beyond the volume boundary, artificially inflating forward projection
values. For an all-ones volume, the GPU returned \texttt{max}=6.14 while the CPU
returned \texttt{max}=2.84 --- a $2.16\times$ mismatch that caused divergence at
every epoch.

The fix is an explicit bounds check in the \texttt{fp\_image} kernel inner loop,
skipping any sample whose fractional voxel index falls outside
$(-0.5,\,\text{dim}-0.5)$:

\begin{lstlisting}[language=C]
if (xi >= 0.f && xi < (float)(Nxz - 1) &&
    yi >= 0.f && yi < (float)(Ny  - 1) &&
    zi >= 0.f && zi < (float)(Nxz - 1)) {
    float4 coord = (float4)(zi+0.5f, yi+0.5f, xi+0.5f, 0.f);
    float density = read_imagef(volume_img, vol_samp, coord).x;
    val += density * dt * rd_norm;
}
\end{lstlisting}

\subsection{Backprojection of All-Ones Scale}

In GPU image mode, \texttt{bp\_ones} was initially computed from a raw all-1.0
image with no preprocessing applied. The resulting scale was off by
$1/\delta_v \approx 313\times$, causing the MLEM denominator to be wrong. The fix
copies a preprocessed (cone-weighted $+$ flip/transpose $/\delta_v$) all-ones buffer
into the image object via \texttt{clEnqueueCopyBufferToImage} before running
backprojection.

% ─────────────────────────────────────────────────────────────────────────────
\section{GPU OpenCL Implementation}
\label{sec:gpu}

\subsection{Mode Overview}

Three GPU modes are implemented, all sharing the same \texttt{proj\_divide},
\texttt{vol\_update}, \texttt{cone\_weight\_hw}, and \texttt{preprocess\_proj}
utility kernels:

\begin{table}[h]
\centering
\begin{tabular}{llll}
\toprule
Mode & Projection storage & Volume storage & Interpolation \\
\midrule
\texttt{gpu-buf} & \texttt{cl\_mem} buffer & \texttt{cl\_mem} buffer & Manual bilinear/trilinear \\
\texttt{gpu-img} & \texttt{image2d\_array\_t} & \texttt{image3d\_t} & Hardware sampler \\
\texttt{gpu-opt} & \texttt{image2d\_array\_t} & \texttt{image3d\_t} & Hardware sampler + LUT \\
\bottomrule
\end{tabular}
\caption{GPU implementation modes}
\end{table}

\subsection{Buffer Mode (\texttt{gpu-buf})}

Work-item mapping: one item per voxel $(ix, iy, iz)$ for backprojection; one item
per detector pixel $(iu, iv, ip)$ for forward projection. Manual bilinear and
trilinear interpolation in the kernel. Global work size: $256^3$ for bp,
$512 \times 512 \times 75$ for fp.

\subsection{Image Mode (\texttt{gpu-img})}

Projections are stored as a \texttt{CL\_MEM\_OBJECT\_IMAGE2D\_ARRAY} with
hardware bilinear sampling in the bp kernel. The volume is stored as a
\texttt{CL\_MEM\_OBJECT\_IMAGE3D} with hardware trilinear sampling in the fp
kernel. This eliminates manual interpolation code and exploits the GPU texture
cache. The image layout uses width $= H$, height $= W$ for the projection array
(matching the preprocessed $[N_\theta][W][H]$ layout) and
width $= N_y$, height $= N_{xz}$, depth $= N_{xz}$ for the volume.

\subsection{Optimized Mode (\texttt{gpu-opt})}

Adds a precomputed float2 cos/sin LUT for the angle loop in backprojection,
replacing one \texttt{cos()} and one \texttt{sin()} call per voxel per angle with
a single coalesced global memory read:

\begin{lstlisting}
float2 cs = angle_cs[ip];   /* preloaded on host: {cos(angle), sin(angle)} */
float ca = cs.x, sa = cs.y;
\end{lstlisting}

For $N_{xz}^2 \times N_y \times N_\theta = 256^2 \times 256 \times 75 \approx 1.26
\times 10^9$ operations per epoch, eliminating transcendental calls reduces bp time
from $\approx 0.05$\,s to $\approx 0.04$\,s (limited by texture bandwidth).

Work-group sizes are tuned to $8\times8\times4 = 256$ threads for bp kernels
(matching the Hawaii GPU wavefront size of 64, occupying 4 wavefronts per group)
and $16\times16\times1 = 256$ threads for fp kernels.

% ─────────────────────────────────────────────────────────────────────────────
\section{Performance Results}
\label{sec:perf}

All timings on AMD Hawaii (Radeon R9 390), $256^3$ volume, $512\times512$ detector,
75 projection angles, 10 epochs.

\begin{table}[h]
\centering
\begin{tabular}{lrrrr}
\toprule
Mode & Time/epoch & Total (10 ep) & Speedup vs CPU \\
\midrule
\texttt{cpu}     & $\sim 7.3$\,s  & 75.9\,s  & $1\times$   \\
\texttt{gpu-buf} & $\sim 1.0$\,s  & 11.9\,s  & $\sim 7\times$  \\
\texttt{gpu-img} & $\sim 0.157$\,s & 1.70\,s & $\sim 47\times$ \\
\texttt{gpu-opt} & $\sim 0.147$\,s & 1.55\,s & $\sim 50\times$ \\
\bottomrule
\end{tabular}
\caption{Reconstruction time, $256^3$ dataset, 10 epochs}
\end{table}

The dominant cost in \texttt{gpu-img} and \texttt{gpu-opt} is forward projection
($\sim 0.12$\,s/epoch), which performs $\sim 2.5\times10^9$ trilinear lookups per
epoch via the image3D hardware sampler. Backprojection costs $\sim 0.04$\,s/epoch.

% ─────────────────────────────────────────────────────────────────────────────
\section{Correctness Validation}
\label{sec:validation}

All GPU modes are validated against the CPU reference by MSE over the full
$256^3 = 16{,}777{,}216$-voxel volume after 10 epochs:

\begin{table}[h]
\centering
\begin{tabular}{lrrrr}
\toprule
Mode & min & max & mean & MSE vs CPU \\
\midrule
\texttt{cpu}     & 0.0000 & 0.9835 & 0.0066 & (reference) \\
\texttt{gpu-buf} & 0.0000 & 0.9838 & 0.0066 & $1.6\times10^{-8}$ \\
\texttt{gpu-img} & 0.0000 & 0.9836 & 0.0066 & $1.9\times10^{-8}$ \\
\texttt{gpu-opt} & 0.0000 & 0.9837 & 0.0066 & $1.9\times10^{-8}$ \\
\bottomrule
\end{tabular}
\caption{Validation results, $256^3$ dataset, 10 epochs}
\end{table}

The residual MSE ($\sim 10^{-8}$) is consistent with float32 rounding differences
between manual trilinear (CPU/buffer) and hardware sampler (image modes). All modes
agree on the reconstructed mean and remain strictly positive, confirming MLEM
convergence is correct.

% ─────────────────────────────────────────────────────────────────────────────
\section{Remaining Work}

The following items are planned for the final report:

\begin{enumerate}
  \item \textbf{512$^3$ dataset}: Run all modes on \texttt{proj\_512\_75.hdf5} and
        report scaled performance and memory usage ($\sim$512\,MB volume).
  \item \textbf{Convergence analysis}: Plot reconstructed volume statistics vs.\
        epoch number for all modes; compare convergence rate GPU vs.\ CPU.
  \item \textbf{Optimization analysis}: Quantify individual contributions of the
        LUT, work-group tuning, and image sampler vs.\ buffer interpolation.
  \item \textbf{Final MSE table}: Full 100-epoch run for all modes with MSE and
        max-diff reported.
\end{enumerate}

% ─────────────────────────────────────────────────────────────────────────────
\section{Summary}

All four reconstruction modes are implemented, validated, and working correctly.
The GPU optimized mode achieves a $\mathbf{50\times}$ speedup over the CPU reference
on the $256^3$ dataset with MSE $< 2\times10^{-8}$. The main technical challenges
were the cone-weight ordering, the fp\_image boundary inflation from
\texttt{CLAMP\_TO\_EDGE}, and the bp\_ones preprocessing scale. All three are
resolved. The final report will cover the 512$^3$ dataset and a full convergence
and optimization analysis.

\end{document}
